\section{Módulo de alarmas}

El módulo recibe señales de alarma del controlador, las emite hacia el entorno y gestiona el ciclo de espera de confirmación para las alarmas críticas.

\subsection{Definición formal}

\begin{align*}
  X &= \{\mathrm{alarmaBaja},\,\mathrm{alarmaMedia},\,\mathrm{alarmaCritica}\}
       \!\times\!\{0\}
     \;\cup\; \{\mathrm{confirmacionEnfermero}\}\!\times\!\{1\} \\[4pt]
  Y &= \{\mathrm{alarmaBaja},\,\mathrm{alarmaMedia},\,\mathrm{alarmaCritica}\}
       \!\times\!\{0\} \\[4pt]
  S &= E_{\mathrm{alarmas}} \times \mathbb{R}_0^+
\end{align*}

El conjunto de estados discretos es:
\[
  E_{\mathrm{alarmas}} = \{
    \mathrm{noAlarma},\;
    \mathrm{alarmaBaja},\;
    \mathrm{alarmaMedia},\;
    \mathrm{alarmaCritica},\;
    \mathrm{espLargo},\;
    \mathrm{espCorto},\;
    \mathrm{repetirCritica}
  \}
\]

El estado general se expresa como $s = (\mathit{est},\,\sigma)$.

\paragraph{Estado inicial.}
\[
  s_0 \;=\; (\mathrm{noAlarma},\;\infty).
\]

\paragraph{Transición externa.}
\[
\dext\bigl((\mathit{est},\,\sigma),\;e,\;(x,\,p)\bigr) =
\begin{cases}
  (\mathrm{alarmaCritica},\;0)
    & p=0,\;x=\mathrm{alarmaCritica} \\[4pt]
  (x,\;0)
    & p=0,\;x\in\{\mathrm{alarmaBaja},\,\mathrm{alarmaMedia}\} \\[4pt]
  (\mathrm{noAlarma},\;\infty)
    & p=1,\;\mathit{est}\in\{\mathrm{espLargo},\,\mathrm{espCorto},\,\mathrm{repetirCritica}\} \\[4pt]
  (\mathit{est},\;\sigma - e)
    & \text{en otro caso}
\end{cases}
\]

\paragraph{Transición interna.}
\[
\dint(\mathit{est},\,\sigma) =
\begin{cases}
  (\mathrm{noAlarma},\;\infty)       & \mathit{est}\in\{\mathrm{alarmaMedia},\,\mathrm{alarmaBaja}\} \\[4pt]
  (\mathrm{espLargo},\;30)           & \mathit{est} = \mathrm{alarmaCritica} \\[4pt]
  (\mathrm{repetirCritica},\;0)      & \mathit{est} = \mathrm{espLargo} \\[4pt]
  (\mathrm{espCorto},\;10)           & \mathit{est} = \mathrm{repetirCritica} \\[4pt]
  (\mathrm{repetirCritica},\;0)      & \mathit{est} = \mathrm{espCorto}
\end{cases}
\]

\paragraph{Función de salida.}
\[
\lambda(\mathit{est},\,\sigma) =
\begin{cases}
  (\mathrm{alarmaCritica},\;0)
    & \mathit{est}\in\{\mathrm{alarmaCritica},\,\mathrm{repetirCritica}\} \\[4pt]
  (\mathit{est},\;0)
    & \mathit{est}\in\{\mathrm{alarmaMedia},\,\mathrm{alarmaBaja}\}
\end{cases}
\]

\paragraph{Avance temporal.}
\[
  \ta(\mathit{est},\,\sigma) \;=\; \sigma.
\]

\subsection{Diagrama de estados}

\begin{figure}[htbp]
\centering
\begin{tikzpicture}[
    node distance=2.2cm and 3.2cm,
    every state/.append style={minimum width=2.4cm}
  ]

  \node[state, initial above] (noA)  {noAlarma};
  \node[state, below left=of noA]  (aB)   {alarmaBaja};
  \node[state, below=of noA]       (aM)   {alarmaMedia};
  \node[state, below right=of noA] (aC)   {alarmaCritica};

  \node[state, below=2cm of aC]    (eL)   {espLargo\\$\sigma{=}30$};
  \node[state, below=2cm of eL]    (rC)   {repetirCritica};
  \node[state, right=of rC]        (eC)   {espCorto\\$\sigma{=}10$};

  \path (noA) edge[bend right=10] node[left]  {\scriptsize alarmaBaja}   (aB)
        (noA) edge                node[right] {\scriptsize alarmaMedia}  (aM)
        (noA) edge[bend left=10] node[right] {\scriptsize alarmaCritica}(aC);

  \path (aB)  edge[bend right=20] node[left]  {\scriptsize $\dint$} (noA)
        (aM)  edge                node[left]  {\scriptsize $\dint$} (noA);

  \path (aC)  edge node[right] {\scriptsize $\dint$: espLargo(30)} (eL)
        (eL)  edge node[right] {\scriptsize $\dint$: repetir}      (rC)
        (rC)  edge node[above] {\scriptsize $\dint$: espCorto(10)} (eC)
        (eC)  edge[bend left=30]
              node[right] {\scriptsize $\dint$: repetir}           (rC);

  \path (eL)  edge[bend left=55, dashed]
              node[above left] {\scriptsize confirm.} (noA)
        (eC)  edge[bend right=60, dashed]
              node[below right] {\scriptsize confirm.} (noA)
        (rC)  edge[bend right=40, dashed]
              node[below] {\scriptsize confirm.} (noA);

\end{tikzpicture}
\caption{Diagrama de estados del Módulo de Alarmas. Las transiciones punteadas representan $\dext$ de confirmación.}
\label{fig:diag_alarmas}
\end{figure}